\chapter{Contexte de l'entreprise}

\section{Pr\'esentation de SEWS Maroc}

\subsection{Le groupe SEWS}

Sumitomo Electric Wiring Systems (SEWS) est une filiale du groupe japonais Sumitomo Electric Industries, l'un des plus grands fabricants mondiaux de c\^ables et harnais automobiles. Fond\'ee en 1917, le groupe Sumitomo Electric emploie plus de 240\,000 collaborateurs dans le monde et pr\'esente un chiffre d'affaires annuel d'environ 30 milliards de dollars \cite{sews2024}.

SEWS Maroc a \'et\'e cr\'e\'e en 2000 \`a K\'enitra, dans la zone industrielle, dans le cadre de la strat\'egie d'expansion internationale du groupe. L'entreprise sp\'ecialise dans la fabrication de harnais \\'electriques pour l'industrie automobile, servant des constructeurs tels que Renault, PSA (Stellantis) et d'autres constructeurs internationaux.

\subsection{SEWS Maroc : chiffres cl\'es}

\begin{table}[htbp]
\centering
\renewcommand{\arraystretch}{1.3}
\begin{tabular}{l l}
\toprule
\textbf{Indicateur} & \textbf{Valeur} \\
\midrule
Date de cr\'eation & 2000 \\
Effectif & 1\,200 employ\'es \\
Chiffre d'affaires & $\approx$ 150 millions EUR \\
Surface d'exploitation & 20\,000 m$^2$ \\
Lignes de production & 15 \\
Capacit\'e annuelle & 3 millions de harnais \\
Principaux clients & Renault, Stellantis, Toyota \\
\bottomrule
\end{tabular}
\caption{Chiffres cl\'es de SEWS Maroc}
\label{tab:sews_chiffres}
\end{table}

\section{Organisation du d\'epartement Maintenance}

\subsection{Structure organisationnelle}

Le d\'epartement Maintenance de \sews est organis\'e sous la responsabilit\'e de M. Mohamed EDDOUCH. Il comprend les \'equipes suivantes :

\begin{itemize}
    \item \textbf{Maintenance \'electrique :} diagnostic et r\'eparation des d\'efauts \\'electriques sur les harnais.
    \item \textbf{Maintenance m\'ecanique :} entretien et r\'eparation des machines et outillages.
    \item \textbf{Maintenance informatique :} support des syst\`emes d'information et des r\'eseaux.
    \item \textbf{M\'etrologie :} calibration et v\'erification des instruments de mesure.
\end{itemize}

\subsection{Le r\^ole du technicien de maintenance}

Le technicien de maintenance joue un r\^ole central dans le bon fonctionnement des lignes de production. Ses missions principales sont :

\begin{enumerate}
    \item \textbf{Diagnostic :} identifier l'origine de la panne en analysant les sympt\^omes.
    \item \textbf{Intervention :} r\'ealiser les r\'eparations n\'ecessaires dans les d\'elais impos\'es.
    \item \textbf{Pr\'evention :} r\'ealiser les op\'erations d'entretien pr\'eventif.
    \item \textbf{Documentation :} enregistrer les interventions et les pi\`eces utilis\'ees.
    \item \textbf{Am\'elioration :} proposer des actions correctives et pr\'eventives.
\end{enumerate}

\section{L'atelier Test \'Electrique}

\subsection{Description de l'atelier}

L'atelier Test \'Electrique est l'un des ateliers critiques de \sews. Il est d\'edi\'e \`a la v\'erification \\'electrique des harnais produits, garantissant leur conformit\'e aux sp\'ecifications des clients. Cet atelier compte 8 lignes de production, chacune \'e\'quip\'ee de machines de test automatis\'ees.

\subsection{Processus de production}

Le processus de production dans l'atelier Test \'Electrique suit les \'etapes suivantes :

\begin{enumerate}
    \item \textbf{R\'eception :} les harnais semi-finis sont r\'eceptionn\'es de l'atelier pr\'e-c\'edent.
    \item \textbf{Pr\'eparation :} les harnais sont positionn\'es sur les gabarits de test.
    \item \textbf{Test \\'electrique :} la machine r\'ealise les mesures de r\'esistance, de court-circuit et d'ouverture.
    \item \textbf{Validation :} les r\'esultats sont compar\'es aux sp\'ecifications du client.
    \item \textbf{Conditionnement :} les harnais valid\'es sont conditionn\'es pour l'exp\'e-dition.
\end{enumerate}

\subsection{Probl\`emes rencontr\'es}

L'analyse du fonctionnement actuel de l'atelier a r\'ev\'el\'e plusieurs probl\`emes :

\begin{itemize}
    \item \textbf{Signalisation manuelle :} les pannes sont signal\'ees oralement ou par radio, sans tra\c{c}abilit\'e num\'erique.
    \item \textbf{Temps de r\'eponse variable :} le d\'elai entre la d\'etection et l'intervention varie de 5 \`a 30 minutes.
    \item \textbf{Absence de statistiques :} aucun outil de suivi des pannes et de leurs dur\'ees n'est disponible.
    \item \textbf{Perte d'informations :} les donn\'ees de production sont partiellement enregistr\'ees manuellement.
    \item \textbf{Difficult\'e d'analyse :} l'identification des tendances et des causes racines est laborieuse.
\end{itemize}

\section{Besoins identifi\'es}

\subsection{Besoins fonctionnels}

L'immersion au sein de l'atelier Test \'Electrique a permis d'identifier les besoins fonctionnels suivants :

\begin{enumerate}
    \item \textbf{Signalisation num\'erique :} remplacement des signaux physiques par des alertes num\'eriques.
    \item \textbf{G\'eolocalisation des pannes :} identification automatique de la ligne et du poste touch\'e.
    \item \textbf{Suivi en temps r\'eel :} affichage en temps r\'eel de l'\`etat de chaque ligne.
    \item \textbf{Historique :} archivage de tous les \'ev\'enements pour analyse ult\'erieure.
    \item \textbf{Notifications :} envoi automatique d'alertes aux \'equipes concern\'ees.
    \item \textbf{Rapports :} g\'en\'eration automatique de rapports et statistiques.
\end{enumerate}

\subsection{Besoins techniques}

\begin{enumerate}
    \item \textbf{Fiabilit\'e :} le syst\`eme doit fonctionner 24h/24, 7j/7, sans interruption.
    \item \textbf{Temps r\'eel :} les donn\'ees doivent \^etre affich\'ees en moins de 2 secondes.
    \item \textbf{Compatibilit\'e :} le syst\`eme doit s'int\'egrer avec l'infrastructure r\'eseau existante.
    \item \textbf{S\'ecurit\'e :} les acc\`es doivent \^etre contr\^ol\'es et les donn\'ees prot\'eg\'ees.
    \item \textbf{Extensibilit\'e :} le syst\`eme doit pouvoir \^etre \'etendu \`a d'autres ateliers.
    \item \textbf{Co\^ut :} le budget allou\'e est limit\'e et le syst\`eme doit \^etre \'economique.
\end{enumerate}

\section{Synth\`ese du chapitre}

Ce chapitre a permis de d\'ecrire l'environnement industriel dans lequel s'inscrit notre travail. \sews, entreprise sp\'ecialis\'ee dans la fabrication de harnais automobiles, dispose d'un atelier Test \'Electrique o\`u les pannes de production sont g\'er\'ees de mani\`ere manuelle et sans tra\c{c}abilit\'e num\'erique. L'analyse des besoins a mis en \'evidence la n\'ecessit\'e d'un syst\`eme Andon num\'erique et connect\'e pour am\'eliorer la r\'eactivit\'e des \'equipes de maintenance et optimiser la gestion des pannes.
